\documentclass[10pt,a4paper]{amsart}
\usepackage[margin=19mm]{geometry}
\usepackage{amsmath,amssymb,mathtools}
\usepackage[scr=rsfs,cal=euler]{mathalfa}
\usepackage{xfrac}
\usepackage[unicode,hidelinks,bookmarks=false]{hyperref}
\newcommand{\ie}{i.e.\ }
\newcommand{\cf}{cf.\ }
\newcommand{\resp}{respectively\ }
\newcommand{\Subsection}{\subsection}
\providecommand{\nobreakdash}{\nobreak\hbox{-}}
\setlength{\emergencystretch}{2em}
\multlinegap0pt
\usepackage{bm}
\usepackage{bbm}
\usepackage{dsfont}
\usepackage{xspace}
\usepackage{cancel}
\usepackage{mathtools}
\usepackage{cleveref}
\usepackage{amsthm}
\makeatletter
\@ifundefined{theorem}{\newtheorem{theorem}{Theorem}}{}
\@ifundefined{lemma}{\newtheorem{lemma}[theorem]{Lemma}}{}
\@ifundefined{proposition}{\newtheorem{proposition}[theorem]{Proposition}}{}
\makeatother
\newcommand{\Aut}{\operatorname{Aut}}
\title[Right groups and the set-theoretic Yang-Baxter equation]{Right groups and the set-theoretic Yang-Baxter equation}
\author{Andrea Albano, Alberto Facchini, Marzia Mazzotta, Paola Stefanelli}
\date{Source: arXiv:2605.25660v1 (CC BY 4.0)}
\begin{document}
\maketitle

\noindent\emph{Source and licence.} Right groups and the set-theoretic Yang-Baxter equation. Version arXiv:2605.25660v1 (CC BY 4.0), distributed under the Creative Commons Attribution 4.0 International licence, as recorded in the arXiv licence metadata for this exact version and shown on the abstract page https://arxiv.org/abs/2605.25660v1. Full rights notice: licenses/arxiv-2605.25660-CC-BY-4.0.txt.\par\smallskip

\noindent\textit{Editorial scope.} Counted unit: the Proposition whose source label is \texttt{None} in the article (source0.tex lines 544--555, source label \texttt{None}), delivered together with the article's own complete original proof of it (source0.tex lines 556--603, one \texttt{proof} environment of 48 lines). No mathematics is written, repaired or re-derived: statement and proof are the article's own text line for line. Required context retained uncounted: EsubsetF (lines 330--332); lambda1 (lines 399--412); somma\_cerchietto (lines 495--507); prop\_GHEF (lines 511--539). Every definition, display and label of those lines is retained unchanged. Omitted from this package and not redistributed: 1--329, 333--398, 413--494, 508--510, 540--543, 604--1485. Nothing inside a retained excerpt is omitted or altered: every retained body is delivered exactly as the article's lines are written, so no omission receipt of any kind accompanies this package. Citations: the retained excerpts carry no citation command at all, so no bibliography entry of the article is transcribed and no reference is redistributed. Reference adaptation: none was needed and none was made; every reference inside the delivered unit resolves inside it, so no unresolved reference or citation is produced. Printed slips kept literal: no printed word, symbol, hypothesis or number of the counted statement or of its proof was altered. Layout and class: the article's own class is \texttt{amsart} with packages including babel, inputenc, amsmath, amssymb, color, enumerate, wasysym, mathrsfs, dsfont, biblatex, xcolor, mathptmx, xy, csquotes; the wrapper is the lane's pinned \texttt{amsart} editorial wrapper at 10pt on A4, which restates the theorem-like environments the retained text needs and adds the article's own macro definitions that the retained text needs (\texttt{\textbackslash Aut}), with extra wrapper packages (bm, bbm, dsfont, xspace, cancel, mathtools, cleveref). The delivered numbering is the unit's own. Assets: no retained span contains a figure environment, a graphics-inclusion command, a PDF-page inclusion, a table or a TikZ picture; no image file is copied into this package, the package declares no asset dependency and no image file is redistributed with it. Licence: version arXiv:2605.25660v1 of this article is distributed under the Creative Commons Attribution 4.0 International licence, as recorded in the arXiv licence metadata for this exact version and shown on the abstract page https://arxiv.org/abs/2605.25660v1. A full rights notice is delivered at licenses/arxiv-2605.25660-CC-BY-4.0.txt. Characters: every retained excerpt is pure ASCII, so the delivered engine, XeLaTeX with its default Latin Modern text font, reports no missing character.
\begin{lemma}\label[lemma]{EsubsetF}
    If $B$ is a left $RG$-semibrace, then every idempotent element in the right group $(B,\circ)$ is also idempotent in the right group $(B,+)$.
\end{lemma}

\begin{proposition}\label[proposition]{lambda1}
    Let $B$ be a left $RG$-semibrace. Define a map $\lambda_a:B \to B$ by setting 
    $$
    \lambda_a(b) = -g_a + a \circ b,
    $$ 
    for all $a,b\in B$.
    Then the following hold:
    \begin{itemize}
        \item[\rm (a)] %for every $a \in B$,  
        $\lambda_a \in \Aut(B,+)$ with inverse given by $\lambda_a^{-1} = \lambda_{h_a^-}$, for every $a\in B$.
           \item[\rm (b)] $\lambda_a(b) = a \circ (h_a^- + b)$, for all $a, b \in B$.
       \item[{\rm(c)}] the map $\lambda:B\to \Aut(B, +)$, $a\mapsto \lambda_a$, is a semigroup homomorphism from $(B, \circ)$ into $\Aut(B, +)$.
     \item[{\rm(d)}] $\lambda_a = \lambda_{h_a}$, for all $a \in B$. 
    \end{itemize}\end{proposition}

\begin{lemma}\label[lemma]{somma_cerchietto}
Let $B$ be a left $RG$-semibrace. Then, for all $a, b \in B$:
    \begin{enumerate}
      \item[\rm (a)] $a \circ b=a+\lambda_a(b)$.
       \item [\rm (b)] $a+b=a \circ \lambda_{h_a^-}(b)$.
    \end{enumerate}
    \begin{proof}
       Let $a, b \in B$. Then we have that $a+\lambda_a(b)=a-g_a+a \circ b=a \circ b.$ Moreover, 
       $$
       a \circ \lambda_{h_a^-}(b) = a \circ h_a^-\circ \left(g_a+b\right)=g_a+b=a+b.
       $$
    \end{proof}
\end{lemma}

\begin{proposition}\label[proposition]{prop_GHEF}
    Let $B$ be a left $RG$-semibrace.
    Then the following hold:
    \begin{itemize}
        \item[{\rm(a)}] $g_a = g_{h_a}$ \ and \ $e_a= \lambda_{a}(f_a)$, for all $a \in B$. In particular, $\lambda_a(0)=e_{h_a}$, for all $a \in B$.
        \item[{\rm(b)}] $\lambda_e\left(H\right) \subseteq H$, for all $e \in E$.
    \end{itemize}
    \begin{proof}
    Let $a \in B$.  By \cref{somma_cerchietto}(a) and \cref{lambda1}(d), we have that
    \begin{align*}
        g_a + e_a = a
        = h_a\circ f_a
        = h_a + \lambda_{h_a}(f_a) 
        = g_{h_a} + \lambda_{a}(f_a). 
    \end{align*}
    Now, $\lambda_{a} \in \Aut(B,+)$ and $\lambda_{a}(f_a)\in E$  by \cref{EsubsetF}. The conclusion follows immediately, thus obtaining the first statement in point {\rm(a)}. 
    In addition,
  $$
\lambda_a(0)=-g_a+g_{h_a}+e_{h_a}=e_{h_a},
  $$ 
  for all $a \in B$.
    Moreover, if we let $e\in E$ and $h\in H$, then %(according to \cref{def2}, we implicitly consider $H = B\circ 0$), 
    by \cref{somma_cerchietto} we deduce that
    \begin{align*}
    \lambda_e(h)=e+\lambda_e(h)=e \circ h=e \circ h \circ 0=\lambda_e(h)\circ 0,
    \end{align*}
    hence $\lambda_e(h)\in H$.
    \end{proof}
\end{proposition}

\begin{proposition}
Let $(B,+,\circ)$ be a left RG-semibrace. Then the following statements hold:
    \begin{enumerate}
        \item[{\rm(a)}] $(E,+,\circ)$ is a left RG-semibrace.
        \item[{\rm(b)}] %$(G + H, +, \circ)$ is  a sub-structure of $B$ \ and \ 
        $(G + H, +, \circ)$ is  a sub-structure of \ $B$ \ and \ $B = (G + H) +  E$.
        \item[{\rm(c)}] $B = H + E= H\circ E$.
        %$(H+ E, +, \circ)$ is  a sub-structure of $B$;
        %\item[{\rm(4)}] $B = H\circ E$.
        \item[{\rm(d)}] $(G\cap H, \circ)$ is a subgroup of the group $(H,\circ)$.
    \end{enumerate}
\end{proposition}

\begin{proof} (a) It is enough to observe that
        \begin{align*}
        e_1\circ e_2 +  e_1\circ e_2
        = e_1\circ e_2 - 0 +  e_1\circ e_2
            = e_1\circ(e_2 + e_2)
            =  e_1\circ e_2,
        \end{align*}
        for all $e_1,e_2\in E$.
        
        (b) Observe that $G+H$ is closed with respect to $+$ by \cref{prop_GHEF}(a). To prove that $G+H$ is closed with respect to $\circ$ it is enough to observe that, by \cref{somma_cerchietto}(a) and again \cref{prop_GHEF}(a),
        \begin{align*}
            b\circ(g+h)
            &= b\circ g - g_{b} 
            + b\circ h\\
            %h_b\circ g - g_{h_b} + h_b\circ h\\
            &= %g_b+g_{\lambda_{h_b}(g)}-g_{h_b} + h_b\circ h 
            g_b+g^{}_{\lambda_{b}(g)}-g_{b} + h_b\circ h
            \ \in \ G + H\,,
        \end{align*}
        for all $b\in B$, $g\in G$, and $h\in H$.
        
(c) Let $b\in B$. Then, by \cref{somma_cerchietto}(a), \cref{lambda1}(d), and \cref{prop_GHEF}(a), we have that
        % $$
        % b = h_b - g_{h_b} + b
        % = h_b - g_b + g_b + e_b
        % = h_b + e_b\in H + E.
        % $$
        \begin{align*}
            b=h_b \circ f_b=h_b+\lambda_{h_b}(f_b)=h_b+\lambda_b(f_b)=h_b+e_b \in H + E.
        \end{align*}
        %\item 
        The second part follows directly from \cref{EsubsetF}.
        
 (d)    Let $x,y\in G\cap H$. Then
\begin{align*}
    x\circ y + 0 
    &= x\circ y - x + x\circ 0 &\mbox{since $x\in G\cap H$}\\
    &= x\circ(y+0)
    = x\circ y&\mbox{since $y\in G$}
\end{align*}
Now, by \cref{prop_GHEF}(a), $\lambda_x(0)=0$, thus,  
by \cref{somma_cerchietto}(b),
\begin{align*}
    x^- + 0 = x^-\circ\lambda_x(0) = x^-\circ 0 = x^-,
\end{align*}
since $x^- \in H$. 
Hence, the claim follows.
\end{proof}

\end{document}
